\documentclass[11pt]{article}
\usepackage{style/blueprintdoc}
\geometry{margin=14mm,top=16mm,bottom=14mm}
\setdoccode{NNV-RC}\setdoctitle{Debate role cards}
\newcommand{\rolecard}[4]{% title, goal, insist, question
\begin{tcolorbox}[enhanced,colback=white,colframe=bpInk,boxrule=0.6pt,arc=0pt,left=2mm,right=2mm,top=1.5mm,bottom=1.5mm,height=6.1cm,fontupper=\footnotesize,borderline={0.4pt}{-1.4mm}{bpRule,dashed},
  overlay={\node[anchor=north east,inner sep=1.2mm,font=\bpmono\tiny,text=bpGraphite] at (frame.north east) {ROLE CARD \quad \faCut\ CUT ALONG THE DASHED LINE};}]
{\bpdisplay\bfseries\normalsize\color{bpInk}#1\par}\vspace{1mm}
\textbf{Your goal.} #2\par\vspace{1mm}
\textbf{You must insist on.} #3\par\vspace{1mm}
\textbf{The uncomfortable question you ask.} \emph{#4}
\end{tcolorbox}}
\newcommand{\decisionsheet}{%
\begin{tcolorbox}[enhanced,colback=bpPaperDark,colframe=bpInk,boxrule=0.6pt,arc=0pt,left=2mm,right=2mm,top=1.5mm,bottom=1.5mm,fontupper=\footnotesize,
  overlay={\node[anchor=north east,inner sep=1.2mm,font=\bpmono\tiny,text=bpGraphite] at (frame.north east) {DECISION SHEET};}]
\begin{tabularx}{\textwidth}{@{}lX@{}}
\textbf{The property, in one sentence} & \rule{0pt}{5mm}\dotfill\\
\textbf{The region} (inputs, ranges, units) & \rule{0pt}{5mm}\dotfill\\
\textbf{$\varepsilon$ / margin} and who chose it & \rule{0pt}{5mm}\dotfill\\
\textbf{Acceptable failure rate} (and for whom) & \rule{0pt}{5mm}\dotfill\\
\textbf{Action on a counterexample} (in development / in production) & \rule{0pt}{5mm}\dotfill\\
\textbf{Who signs}, who is accountable when the specification was wrong & \rule{0pt}{5mm}\dotfill\\
\textbf{Out of scope for verification} and how it is addressed instead & \rule{0pt}{5mm}\dotfill\\
\end{tabularx}
\end{tcolorbox}}
\begin{document}
% ------------------------------------------------------------------ scenario 1
\bptitleblock{Scenario 1 · Medical triage}{A neural network turns vital signs into an urgency class}{DOCUMENT NNV-RC \quad|\quad REV A \quad|\quad 16 September 2026 \quad|\quad SCENARIO 1 OF 3}
\begin{notebox}\footnotesize
\textbf{Facts on the table.} A hospital trust wants to replace its rule-based early-warning score with a learned model. In this afternoon's exercise the 96\,\%-accurate model under-triaged a patient with SpO\textsubscript{2} 86\,\% (found in 0.5\,s); the corrected model still failed at exactly 90\,mmHg; the model trained with the properties in the loop proved all five clinical limits but over-triages at the thresholds (92\,\% agreement with the strict rule). Elsewhere: a certified neural dosing controller for vancomycin was proved never to exceed the safe concentration over an unbounded sequence of doses — under a simplified pharmacokinetic model and synthetic training data (Sirman et al., SAIV 2026). An analysis reported in \emph{Biomedical Safety \& Standards} (Feb 2026) attributes nearly half of recent FDA medical-device recalls to AI-enabled devices, mostly for diagnostic or measurement error. Under Article~15 of the EU AI Act (obligations from 2 August 2026) high-risk systems must reach ``an appropriate level of accuracy, robustness, and cybersecurity''.
\end{notebox}
\begin{tabularx}{\textwidth}{@{}XX@{}}
\rolecard{Regulator}{Approve only what can be shown safe for every patient the device will see, and be able to explain the approval later.}{Written properties with named regions; evidence for \emph{all} inputs in scope, not a test set; a plan for what happens when a counterexample is found after deployment.}{Which clinical rule did you decide not to prove, and why?} &
\rolecard{Vendor}{Ship a model that is accurate, fast, and defensible; keep the guarantee without losing the accuracy that sold it.}{A clear statement of the operating region so that out-of-region inputs are not counted as failures; margins that do not make the product useless.}{Do you want the 96\,\% model with a hole or the 92\,\% model with a proof? You cannot have both today.}\\[4mm]
\rolecard{Clinician / patient}{Never miss the deteriorating patient; not be drowned in false alarms either.}{Over-triage has a cost (beds, attention, alarm fatigue) that must be on the sheet; a human can always override; the thresholds are clinical decisions, not engineering ones.}{Who told the model that 92\,\% saturation is ``fine''?} &
\rolecard{Verification engineer}{Say precisely what has and has not been proved, and refuse to let a proof be quoted beyond its region.}{Inputs must be bounded; flags pinned; every threshold a unit away from where the labels change; every counterexample reproduced and root-caused.}{The proof says nothing about a patient whose readings are wrong. What does?}\\
\end{tabularx}
\vspace{2mm}\decisionsheet
\newpage
% ------------------------------------------------------------------ scenario 2
\bptitleblock{Scenario 2 · Authentication}{Live facial recognition decides who gets stopped}{DOCUMENT NNV-RC \quad|\quad REV A \quad|\quad 16 September 2026 \quad|\quad SCENARIO 2 OF 3}
\begin{notebox}\footnotesize
\textbf{Facts on the table.} According to figures reported by PA in 2026, the Metropolitan Police used live facial recognition 231 times in 2025 and scanned about 4.2\,million faces; ten people were falsely alerted, eight of them Black; none of the ten was arrested. Shaun Thompson, an anti-knife-crime youth worker, was stopped and detained on a false match in 2024 — his passport and bank cards were ``not enough''; on 21 April 2026 the High Court dismissed his and Big Brother Watch's challenge and ruled the deployment lawful (The Register, 22 April 2026). Industry reports in 2026 describe deepfake camera injection and voice cloning sold as a service against biometric onboarding (vendor sources; treat the figures with care). Robustness of a face-matching network can be stated as a property around known inputs; ``never misidentifies a person'' cannot.
\end{notebox}
\begin{tabularx}{\textwidth}{@{}XX@{}}
\rolecard{Regulator}{Permit only deployments whose error behaviour is specified, measured and reviewable, with equal protection across groups.}{A stated region (lighting, pose, resolution) and $\varepsilon$; error rates \emph{per group}, not overall; what a false alert leads to, in minutes and in law.}{Ten in four million sounds small. Ten of whom?} &
\rolecard{Vendor}{Deliver a system that catches wanted people and passes an audit.}{The operating envelope the numbers were measured in; the threshold that trades false alerts against misses is set by the customer, not the vendor.}{If you demand zero false alerts, you are asking for a system that alerts on nobody. Which do you want?}\\[4mm]
\rolecard{Affected person}{Not be stopped, filmed and doubted because of a similarity score.}{A counterexample here is a person; the ``action on a counterexample'' must be a five-minute check, never an arrest on the score alone; a way to contest the match.}{What would have convinced the officer that the machine was wrong?} &
\rolecard{Verification engineer}{State exactly what robustness around a reference image does and does not cover.}{An adversary who \emph{chooses} the perturbation (glasses, make-up, injection) is stronger than random noise; the property must say which adversary it assumes.}{Robust to what, chosen by whom?}\\
\end{tabularx}
\vspace{2mm}\decisionsheet
\newpage
% ------------------------------------------------------------------ scenario 3
\bptitleblock{Scenario 3 · Autonomous control}{A safety envelope you can prove versus a perception system you cannot}{DOCUMENT NNV-RC \quad|\quad REV A \quad|\quad 16 September 2026 \quad|\quad SCENARIO 3 OF 3}
\begin{notebox}\footnotesize
\textbf{Facts on the table.} On 20 April 2026 an empty Waymo robotaxi drove into a flooded road in San Antonio and was swept into a creek; the fleet's software was written to \emph{slow, but not stop} at water it judged ``potentially untraversable''; 3,791 vehicles were recalled and fixed over the air (CNBC, 12 May 2026). In March 2026 the US regulator upgraded its investigation of Tesla's Full Self-Driving to an engineering analysis over crashes in reduced visibility, covering about 3.2\,million vehicles. In the tutorial you saw a neural controller \emph{proved} to keep a car on the road under crosswind and sensor noise (Vehicle + Marabou, with a machine-checked proof of the whole system in Agda), and the ACAS Xu collision-avoidance network \emph{disproved} for Property 3 in under a second.
\end{notebox}
\begin{tabularx}{\textwidth}{@{}XX@{}}
\rolecard{Regulator}{Approve a system whose safety case rests on evidence about every situation in the operating domain, and know who to hold responsible.}{A written operating domain; proofs for the control layer; an honest account of the perception layer that cannot be proved; incident reporting on time.}{Whose signature is on the sentence ``slow, but do not stop''?} &
\rolecard{Vendor}{Deploy widely, fix fast, keep the fleet running.}{Verification of the controller is affordable and should be in the pipeline; perception is monitored statistically; over-the-air fixes are part of the safety case, not an admission.}{A recall in ten days is the system working. What would you have preferred?}\\[4mm]
\rolecard{Road user}{Not be the counterexample.}{The difference between ``the car never leaves the road'' (a proof about a controller) and ``the car never hits a child'' (not a property of pixels) must be explained in plain words before deployment; a human must be able to stop the vehicle.}{Which of the two sentences are you actually promising me?} &
\rolecard{Verification engineer}{Prove what can be proved and label the rest.}{Controllers and envelopes: prove them, in the units of the problem. Perception: bound it with robustness properties around known scenes and say so. Specifications are reviewed like code — they were the bug in San Antonio.}{The proof of the controller assumes the sensor is right. Who owns that assumption?}\\
\end{tabularx}
\vspace{2mm}\decisionsheet
\end{document}
